\chapter{Założenia projektowe}

Głównym założeniem stawianym całemu projektowi jest wykonanie prostego i 
wygodnego w obsłudze środowiska, pozwalającego na wykonywanie badań i
eksperymentów związanych z omawianą tematyką. Poprzez pojęcie środowiska, 
autor ma na myśli nie tylko platformę sprzętową, ale również strukturę
oprogramowania. Sprzęt ma być prosty w obsłudze i możliwe odporny 
na błędy użytkownika, a struktura oprogramowania powinna udostępniać
zestaw narzędzi pozwalających na łatwy proces prototypowania, symulacji, 
komplilacji i ostatecznie wdrożenia zaimplementowanego rozwiązania.

\par Autorowi pracy zależy, żeby wykonany projekt miał jak najszersze 
horyzonty edukacyjne. W związku z tym, platforma sprzętowa, poza wykorzystaną
jednostką obliczeniową, zostanie zaprojektowana i sporządzona samodzielnie w 
postaci płytki drukowanej. Oprogramowanie napisane na docelową 
platformę nie będzie korzystało z bibliotek udostępniających gotowe funkcje
realizujące algorytmy będącymi znaczącymi elementami pracy.
Wykorzystanie zewnętrznych bibliotek zostanie zamknięte na poziomie 
prototypowania i symulacji.

\par Dodatkowo postanowiono, że wykorzystane narzędzie wspomagające
proces projektowania sprzętu jak i implementacji oprogramowania będą 
na licencjach wolnego oprogramowania.

\section{Wymagania stawiane części sprzętowej}

Z ogólnych zadań zaprojektowanego i sporządzonego układu elektronicznego 
można wyróżnić:
\begin{enumerate}

    \item Obsłużenie i odpowiednie przetworzenie wejściowych sygnałów analogowych 
    z mikrofonów pomiarowych pod kowersję do domeny cyfrowej z użyciem 
    osadzonych na płytce przetworników analogowo-cyfrowych.

    \item Obsłużenie i odpowiednie przetworzenie wyjściowych sygnałów analogowych
    pochodzących z osadzonych na płytce przetworników cyfrowo-analogowych pod
    poprawne dostarczenie ich na głośniki.

    \item Zapewnienie działającej komunikacji z wybraną platformą obliczeniową 
    poprzez wybrany protokół cyfrowy, przez który przesyłane będą przekonwertowane
    z użyciem osadzonych na płytce przetworników analogowo-cyfrowych 
    do domeny cyfrowej wejściowe, pochodzące z mikrofonów, sygnały analogowe, 
    oraz końcowo trafiające na głośniki dane wyjściowe do przetworników 
    cyfrowo-analogowych.

    \item Zapewnienie odpowiednich zasilań dla sekcji analogowych oraz cyfrowych.
    Możliwość zasilania wybranej platformy obliczeniowej z poziomu płytki.

\end{enumerate}

Aby powyższe zadania wykonać, układ musi spełnić następujące wymagania:
\begin{enumerate}

    \item Dla wejściowych sekcji analogowych: 
    \begin{enumerate}

        \item Konwersję symetrycznego sygnału z mikrofonów 
        do odpowiednio wzmocnionej sumy.

        \item Filtrację dolnoprzepustową ograniczającą pasmo sygnału wejściowego
        do mieszczącego się w zakresie od 50 do 500 Hz.

        \item Odpowiednie wyskalowanie przetworzonego sygnału, aby mieścił się
        w pełnym zakresie dynamicznym przetwornika, jak i nie przekraczał
        dopuszczalnej granicy.

    \end{enumerate}

    \item Dla wyjściowych sekcji analogowych:
    \begin{enumerate}

        \item Odpowiednie przetworzenie sygnału wyjściowego przetworników
        cyfrowo-analogowych, opierające się na skalowaniu i filtracji
        rekonstrukcyjnej.

        \item Przekształcenie wyjściowego sygnału na symetryczną parę,
        zdefiniowaną przez standard złącza XLR.

    \end{enumerate}

    \item Dla wybranego protokołu cyfrowego:
    \begin{enumerate}
        
        \item Synchronizację wszystkich przetworników poprzez wystawienie
        wspólnego sygnału zegarowego o konfigurowalnej częstotliwości.

        \item Jednoczesną obsługę dwóch kanałów wejściowych (RX) i 
        dwóch kanałów wyjściowych (TX). Oznacza to, że protokół 
        musi wspierać multipleksowanie kanałów, jeżeli wykorzystane
        zostaną pojedyncze linie danych na odpowiednim kierunku transmisji,
        lub odpowiednią konfigurację selektorów (CS), ewentulanie selektora (FS), 
        jeżeli zostaną wykorzystane pojedyncze linie danych na kanał.

    \end{enumerate}

    \item Dla wybranych przetworników:
    \begin{enumerate}

        \item Relatywnie wysoka rozdzielczość, wynosząca minimalnie 12 bitów.
        
        \item Wsparcie dla identycznego protokołu cyfrowego.

    \end{enumerate}

    \item Dla sekcji zasilania:
    \begin{enumerate}

        \item Generację wszystkich potrzebnych zasilań na podstawie jednego,
        dodatniego doprowadzonego z zewnątrz.

        \item Generację symetrycznych zasilań dla wzmacniaczy pracujących
        w sekcjach analogowych o względnie wysokiej dokładności.

        \item W przypadku wyboru mikrofonów pojemnościowych, generację
        odpowiedniego zasilania phantomowego, typowo mieszczącego się w zakresie
        od +12 do +15 V.

        \item Generację zasilań dla przetworników i zewnętrznej platformy
        obliczeniowej.

        \item Generację dokładnych napięć referencyjnych dla przetworników
        (jeżeli będą potrzebne).

    \end{enumerate}

\end{enumerate}

W przypadku urządzeń zewnętrznych wymagania są następujące:
\begin{enumerate}
    
    \item Dla wykorzystanych mikrofonów:
    \begin{enumerate}

        \item Stosunkowo niska cena i wysoka dostępność.

        \item Typ - standardem w omawianej tematyce są mikrofony \textbf{pojemnościowe}.
        Znane są z wysokiej czułości i płaskiej charakterystki częstotliwościowej, jednak
        wymagają dodatkowego zasilania, zwanego zasilaniem phantom.
    
        \item Charakterystyka kierunkowa - superkardoidalne lub hiperkardoidalne 
        mogą lepiej skupić się na kierunkowym źródle dźwięku, lecz w przypadku hałasu 
        w otwartej przestrzenii odbierany dźwięk będzie słabszy, jeżeli jego źródło 
        znajdzie się poza docelowym obszarem działania. Tematyka 
        projektu skupia się na hałasie w przestrzeni, a źródła dźwięku ustawione są
        względem siebie pod pewnym kątem, więc docelową charakterystyką jest 
        \textbf{charakterystyka kulista}.
    
        \item Odpowiedź częstotliwościowa - ideowo \textbf{płaska}. Zakres nie odgrywa
        znaczącej roli, bo analizowane sygnały są wolnozmienne (pasmo od 50 do 500 Hz, 
        gdzie typowe zakresy posiadają górną granicę definiowaną w dziesiątkach kiloherców).
    
        \item Czułość - informuje o wartości napięcia dla danego ciśnienia aktustycznego.
        Wynika z następującej zależności:
        \begin{equation} 
                S_{\text{dBV}} = 20 \text{log}_{10} \frac{S_{\text{mV/Pa}}}{U_{\text{ref}}}
        \end{equation}
        gdzie $U_{ref} = 1000 \text{ }\frac{\text{mV}}{\text{Pa}}$. 

    \end{enumerate}

    \item Dla wykorzystanych zestawów głośnikowych:
    \begin{enumerate}

        \item Z racji, że zrezygnowano z projektowania końcówki mocy, wymagany 
        jest wbudowany wzmacniacz. Innymi słowy, zestawy muszą być \textbf{aktywne}.
    
        \item Analizowane sygnały są wolnozmienne (pasmo 100 do 500 Hz), więc 
        wymagane jest wsparcie dla tonów niskich. Innymi słowy na wyposażeniu zestawu
        musi być głośnik niskotonowy (ang. \textit{woofer}), zdolny do emisji
        sygnałów akustycznych o niskiej częstotliwości.
    
        \item Z uwagi na zakres poziomu akustycznego analizowanych sygnałów, 
        zestaw musi być zdolny do osiągnięcia jego górnego limitu, czyli użytecznego 
        90 dB SPL.

    \end{enumerate}
\end{enumerate}


\section{Wymagania stawiane części programowej}
